\chapter[Facts, Explanations, and Values]{Can Facts, Explanations, and Values Be Separated?}
\section{A Clod of Earth Beside a Skull}
First, pick up something inconspicuous: a piece of earth.

Not ordinary garden soil, but a yellowish clod in a glass case in a museum, containing particles almost too small to see. It comes from a cave in the Zagros Mountains of northern Iraq and dates to about fifty to seventy thousand years ago. The label identifies the particles as pollen, powder released by ancient flowering plants.

Why look at soil? Because beside it lies a human skeleton, more precisely a Neanderthal skeleton. Neanderthals were our close relatives, ancient humans living in Europe and western Asia who disappeared about forty thousand years ago. When this individual was buried, the soil beneath the body contained large quantities of pollen, not scattered grains but clusters.

Where did the pollen come from?

If the wind brought it, this is ordinary soil. But the archaeologist directing the excavation proposed an answer that shook the field: the pollen came from bunches of flowers placed in the grave by the dead person's relatives. More than fifty thousand years ago, people we had long called cave-dwelling savages may have held a flower-adorned funeral for a companion.

A clod suddenly opens onto a huge question: did they possess humanity? What entitles us to ask? Behind it lies this chapter's larger question. Pollen is there: a fact. It was a funeral: an explanation. They therefore had humanity and deserve respect: a value judgment. Can these three be separated completely? Should they be?

\section{An Excavation in the Zagros Mountains}
Now enter a scene with me.

Time: a summer in the late 1950s. Place: Shanidar, an immense cave in the Zagros Mountains of Iraqi Kurdistan. Its entrance faces north. Inside it is cool and smells of limestone and bat droppings. A dry river valley lies below; distant shepherds drive sheep, their bells sounding intermittently.

Dozens of archaeologists and locally hired Kurdish workers occupy the cave. The American archaeologist Ralph Solecki has worked here for several seasons. Excavation is tedious: small trowels and brushes peel away millennia of soil, layer by layer. Each layer is sieved; bone fragments, stone tools, and charcoal are numbered and bagged. There is no electricity. At night they light fires at the entrance, plagued by mosquitoes.

Then a trowel deep inside meets bone. Not one bone, but skeleton after skeleton, several Neanderthals. Workers lie close to the ground, switching to bamboo picks and brushes the size of writing brushes to remove soil from the gaps. This is excavation's most demanding moment. Any brushstroke could destroy information fifty thousand years old.

Beside the skeleton labeled Shanidar 4, Solecki notices soil of a different color. He bags surrounding samples and sends them to a laboratory. Years later, the French pollen specialist Arlette Leroi-Gourhan counts them under a microscope. The pollen is unusually abundant and clustered, much of it from brightly flowering plants.

Solecki concludes that mourners placed flowers around the body, perhaps individual stems or woven garlands. He later makes the image the title of a book, Shanidar: The First Flower People. The conclusion is deeply moving. Death, mourning, ritual, things we imagined belonged only to modern humans, existed in another kind of human too. Newspapers compete to report it. The Neanderthal flower burial becomes one of twentieth-century archaeology's most famous stories.

Remember this cave. The questions ahead grow from those bags of earth.

\section{What Did the Pollen Actually Say?}
Let us spread out the difficulties without rushing to conclude.

The pollen really is there. A microscope does not lie, at least not deliberately. But between pollen being there and relatives offering flowers lies a wide river.

Later researchers propose other crossings. A gerbil inhabiting Shanidar stores plants in its burrows. Pollen on plants carried inside could fall into the soil in clusters. Others suggest bees and other insects might bring pollen into caves. On these accounts, the pollen records a rodent's storeroom or an insect passage, not a funeral.

Has the flower burial therefore been disproved? Not necessarily. Its supporters argue that pollen quantity, distribution, and relation to the skeleton may not be fully explained by rodents. Recent excavations have also found remains that may have been deliberately placed. The debate remains unresolved.

Notice that what we have just done belongs to three different levels:

\begin{itemize}
\item \textbf{What is there?} Clusters of pollen in soil beside a skeleton. This is the factual level: in principle, anyone counting through a microscope should find the same thing.
\item \textbf{Why is it there?} Did a funeral leave it, or gerbils bring it? This is explanation. Competing accounts try to make sense of one fact, and more evidence must decide between them.
\item \textbf{What should follow?} If this really was a funeral, should textbooks change to acknowledge Neanderthals' humanity and grant them the respect we give our ancestors? This is value. Microscopes cannot settle it.
\end{itemize}
Trouble arises when the levels are quietly glued together. Someone jumps from pollen to mourning, then to a demand for renewed respect. Another, unwilling to share humanity with Neanderthals, tries to diminish even the clustered-pollen fact. Emotion runs ahead of evidence; positions choose their facts.

The problem is neither ancient nor confined to archaeology. You encounter it daily. Your desk mate says, ``The teacher kept us ten minutes late again'' (fact?), ``He just dislikes our class'' (explanation), ``He doesn't deserve to teach'' (value). Said in one breath, they sound like one thing but are three. Arguments often arise because nobody separates them.

Before reading on, take a position in the cave. If gerbils ultimately prove responsible for the pollen, would Neanderthals fall in your estimation? Should they?

\section{Who May Tell the Past?}
The Shanidar dispute seems calm because its subjects died fifty thousand years ago. Nobody will burst into tears before us. Change the participants and the same question becomes painful.

Consider the Americas. In 1492, Columbus's fleet reached the Caribbean. For many years, the United States commemorated the event with Columbus Day in October. Textbooks portrayed a heroic discoverer, and schoolchildren acted out his voyage. Since the late twentieth century, more states and cities have renamed the occasion Indigenous Peoples' Day. Another community supplies the reasons. For Indigenous Americans, 1492 did not begin their discovery. They already had towns, fields, languages, and gods and needed nobody to discover them. The arrival brought disease, enslavement, and loss of land. One history appears as a festival on one calendar and a day of mourning on another.

Now give the less favored side its full case. We practice this twice in the chapter, beginning with defenders of the commemoration.

Calls for renaming are loud, but the defenders deserve a hearing, not ridicule. Their reasons have three layers. First, they commemorate not only Columbus but family stories. The American holiday was closely connected with Italian immigrants. For families whose ancestors crossed the ocean and suffered discrimination, it affirmed that they belonged to the country. Removing it sounds like withdrawing that affirmation. Second, they argue that past people should not be judged solely by today's standards. Nearly all fifteenth-century voyagers participated in acts now considered crimes; that standard would leave no historical holidays. Third, they fear calendar revision becoming a competition in correct public postures while practical improvements for Indigenous communities, such as reservation education and health care, are neglected.

You may reject every argument, but first acknowledge the real questions: who is doing the remembering, which standards judge the past, and how symbols relate to tangible benefits. Renaming's supporters have an equally full case. Calendars are not neutral. Printing a person's name on a national holiday tells children that this is someone we honor together. One group's festival should not rest on another's wounds. And the defense that we cannot use today's standards forgets contemporary protesters using contemporary standards. The Spanish missionary Bartolomé de las Casas recorded atrocities he witnessed in Columbus's era and spent his life advocating for Indigenous people. There was never only one standard then, either.

Now an older Chinese scene shows ancestors arguing just as skillfully. The Sui Grand Canal required incalculable labor, many died building it, and the dynasty soon fell amid rebellion. How should history describe the canal? Tang people already disagreed. Many said it destroyed the Sui. Yet the late Tang poet Pi Rixiu wrote in ``Meditating on the Past by the Bian River'':

\begin{quote}
All say this river brought the Sui's fall; still its waters serve a thousand li. Without those floating palaces and dragon boats, its merit might stand alongside Yu's.
\end{quote}
Roughly: everyone blames the canal for the dynasty's fall, yet long-distance shipping still depends on it. Set aside Emperor Yang's pleasure cruises and the achievement could rival Yu the Great's. Notice the move. Pi does not erase forced labor and deaths; the floating palaces and dragon boats remain on the charge sheet. Nor does he deny the canal's later benefits. He separates what kind of ruler Emperor Yang was from what kind of engineering project the canal was, and answers separately. A poet eleven centuries ago already demonstrated this chapter's skill.

Thus disputes over who may tell the past often concern more than facts. Ships arrived in 1492 and the canal was dug; neither can be denied. The disputes concern which fact leads, whose standard evaluates, and whose voice counts. Fact, explanation, and value wind together through every history book. Our task is not pretending them absent, but finding each thread.

\section{Thinking Bridge: Three Baskets}
Finding the threads can become a portable tool. Whenever you hear a weighty statement, first ask which basket it belongs in.

\textbf{Basket one: fact. What is the case?}

It asks what the world is like. Is there pollen in the soil? How many kilometers is the canal? Did ships arrive in 1492? What is today's temperature? Such questions have right and wrong answers decided by evidence that can be measured, checked, and rechecked. If you report pollen clusters, I can count your soil sample again. If the result differs, you are wrong.

\textbf{Basket two: explanation. Why?}

Why is the pollen there? Why did Sui fall? Why did the teacher run late? These questions are slipperier than factual ones. The same fact often fits conflicting accounts, like funeral and gerbils. Evidence still decides, but an extra step asks which account explains more and leaves fewer puzzles. Can gerbils explain the pollen's location? Can the moving funeral account explain other soil layers? Good explanations submit to questioning. Bad ones merely repeat themselves.

\textbf{Basket three: value. What ought to be?}

It asks good or bad, right or wrong, should or should not. Should Columbus be commemorated? Do Neanderthals deserve respect? Is keeping a class late right? No microscope rules here. Shared standards, fairness, compassion, responsibility, promises, and arguments about them are the judges.

The tool's whole point is \textbf{sort first, then work}. Each basket has its own rules. Put the wrong referee in charge and the game becomes impossible to follow.

Consider common wrong-referee accidents:

\begin{itemize}
\item Using factual measures to judge values: ``He's hopeless. He scored only 40 in math.'' Forty is a fact; hopeless is a value verdict with a vast argumentative gap between them. One test cannot support a final judgment on a person.
\item Using value wishes to rewrite facts: ``That claim is too hurtful to be true.'' Being painful and being false are different. Truth is not responsible for our comfort.
\item Disguising explanation as moral acquittal: ``He hits people because he was abused as a child, so he's done nothing wrong.'' Explaining causes leaves the action's account unsettled. Understanding the road to a cliff does not make jumping right.
\end{itemize}
Conversely, sorting loosens many knots. Two people argue over whether global warming is a hoax until their friendship breaks. Examine the dispute: one concerns temperature data, the other whether emissions reduction should require lifestyle changes. Their views of the data may barely differ. Next time a sentence provokes you, ask which basket it currently occupies before replying.

\section{A Self-Evident Declaration and Its Crack}
Now a short primary source, perhaps history's most famous act of putting values into facts' pocket.

In 1776, representatives of thirteen North American colonies signed the Declaration of Independence. Near its opening appears the sentence heard around the world:

\begin{quote}
We hold these truths to be self-evident, that all men are created equal, that they are endowed by their Creator with certain unalienable Rights, that among these are Life, Liberty and the pursuit of Happiness. ---United States Declaration of Independence, 1776
\end{quote}
Sort the sentence. All are created equal: is it fact? Look around. Some are tall, others short; rich or poor; healthy from birth or born ill. As factual description, it meets obstacles everywhere. It is a declaration of value, not what people \textbf{are like}, but \textbf{how they should be treated}. In dignity and rights, nobody should be inherently inferior. The drafters understood this and supplied a clever device: we hold, self-evident. They did not claim everybody could see it, since reality contradicted it. They said this value would be their starting point, an uncompromised premise by which to judge reality.

Now the famous crack. Thomas Jefferson, the principal drafter, enslaved hundreds of Black people over his life. The words of equality came from someone who denied equality to those he enslaved. Contemporaries and later readers seized on this: fine words, hypocrisy.

Wait. Here lies one of the chapter's easiest misjudgments. Walk slowly.

In 1852, the Black abolitionist Frederick Douglass, himself escaped from slavery, was invited to speak at an Independence Day celebration. His point was roughly: what does this country's holiday mean to the enslaved? The freedom celebrated is yours, not ours. He did not say the Declaration was false and should be discarded. Instead he used it as a weapon: fulfill the self-evident promise you yourselves wrote. More than eighty years later, Lincoln's Gettysburg Address again invoked a nation founded eighty-seven years earlier, conceived in liberty and dedicated to human equality, treating the statement as a debt to be paid.

So declaring the Declaration hypocritical is too quick. A fuller picture shows a moral promise written by someone who could not fulfill it, signed by a country that could not fulfill it, then held aloft for two centuries by successive generations of the mistreated as evidence of a debt. Value language deceives and compels. It is a fig leaf and an IOU. See only the fig leaf and you miss its power to change the world. See only the IOU and you forget how long payment was delayed. Fact, explanation, and value knot together here, and that knot makes history worth reading.

\section{One Ocean, Three Explanations}
Now a demanding exercise: confront one grave fact with different explanations and examine each one's strength and boundary.

First establish the factual level. From the sixteenth through nineteenth centuries, the European-run Atlantic slave trade transported shipload after shipload of Africans to the Americas. Surviving manifests, customs records, and port accounts let historians count more than ten million Africans forced onto ships and carried across the Atlantic, with most estimates near twelve million, alongside vast numbers who did not arrive alive. These are inspectable accounts, not legends. The trafficked people are not silent numbers. Archives preserve names and resistance. Revolts repeatedly broke out aboard ships; escapees established free communities in the American interior. They always struggled and acted within their own destinies, though deprived of the possibility of winning.

Why did this happen? Three explanations have their own territories.

\textbf{Explanation one: money.} American sugar plantations and cotton fields demanded enormous labor. Indigenous people died in large numbers from disease, European immigrants were unwilling to do the work, and enslaved Africans became a profitable business. Profits returned to Europe and enriched ports, banks, and factories. This account explains scale, routes, and duration powerfully: nobody maintains an unprofitable business for three centuries. Its limit is clear too. However carefully calculated, the business first required a justification for treating people as cargo. Interest alone cannot explain why European merchants did to Africans what they would not do to their compatriots.

\textbf{Explanation two: racism.} Europeans developed accounts of Africans as inferior and destined for servitude, making enslavement seem natural. This fills money's gap by explaining how moral barriers were removed. Its limit is chronological. Many systematic racist theories were elaborated \textbf{after} the trade already operated on a large scale. They resemble licenses issued retroactively to a working business more than its sole engine. Ideas kill, certainly, but interests often summon the ideas first.

\textbf{Explanation three: power and institutional openings.} Europeans did not simply enter Africa alone to capture people. Ships largely remained on the coast, supplied by local networks of warfare, raids, and trade. Some African kingdoms and merchants participated deeply, exchanging captives for firearms and goods. This account sees structure. Slavery was already familiar in Africa and Eurasia; Atlantic demand magnified an old institution into a monster. Evil often operates not through a conspiracy of demons but through countless people calculating their own interests within established systems. Its limit is that assigning responsibility to structure can dilute concrete guilt. European merchants, slave-ship owners, and plantation owners were major beneficiaries and drivers. They held the ledgers. That must remain clear.

All three capture part of the truth. Notice again the distinction this chapter keeps making: \textbf{explaining why something happened does not excuse it}. Calculating slave-ship profits does not absolve owners. Precisely because mechanisms become clear, we can identify where to stop them. Understanding how someone became a murderer does not mean asking a court to release them. Sympathetic understanding, entering an actor's situation to see available knowledge and choices, and critical judgment, standing with the harmed to identify responsibility, need one another. Criticism without understanding fires blanks; understanding without criticism evades judgment.

Remember too that condemning the slave trade is a value judgment, but it \textbf{does not and need not} alter any fact. Manifest figures neither grow with our anger nor shrink with someone else's defense. This is the healthiest relationship among fact, explanation, and value.

\section{Further Challenge: No ``Ought'' from ``Is''}
Now the chapter's central theoretical discovery, from the eighteenth-century Scottish philosopher David Hume. In A Treatise of Human Nature, he observed a previously unremarked shift. Moral discussions begin with what the world is like, divine arrangements, human nature, facts, then quietly switch to what ought to happen: obedience, restraint, kindness. Authors never explain the shift, yet these are different kinds of statement. An ought cannot simply be inferred from an is.

Later readers called this Hume's guillotine, separating fact and value: however many factual propositions you collect, they do not yield a value conclusion unless a value premise is added.

See what the blade cuts in ordinary life. ``Humans are naturally selfish, so fair distribution is a dream.'' Even if natural selfishness were fact, which is disputed, it would not establish that fairness should not be pursued without an unstated premise: whatever is natural should be left alone. Yet we never consistently follow that rule. People naturally become ill, and we treat them. Children are born unable to read, and we teach them. Bring the hidden premise into daylight and it fails.

A darker example says, ``Research shows children in that region average lower grades, so investing more resources in them is wasteful.'' Even if the data are true, not worth investing is a value judgment carrying another premise: people's worth depends on average grades. Expose that premise and many speakers would blush. Hume's guillotine does not prohibit values. It forces value premises into public view for debate. Whenever you hear facts are this way, therefore things ought to be this way, feel between the clauses for the smuggled premise.

Follow Hume a step farther and you meet our title's central difficulty. If values cannot be deduced or settled that way, can history and sociology, which study people every day, still be objective?

Max Weber addressed this in the early twentieth century, with an answer in two halves. First, values \textbf{do} enter research and cannot be excluded, because they decide \textbf{what we study}. A scholar concerned with equality may inspect slave-trade accounts; one concerned with order, imperial law codes. Choosing a subject is never neutral, and pretending otherwise only conceals one's position. Weber considered this no disgrace, but something to acknowledge. Second, after selecting the question, values should step aside. Research must obey another discipline: disclose evidence's origins, retain unwelcome material, expose every inferential step to checking and refutation, and publish conclusions even when they disappoint you.

Together, the halves yield this chapter's key statement: \textbf{objectivity is not having no position, but submitting methods, evidence, and reasoning to public examination}. Positionless people do not exist. There are those who hide positions and let them secretly shape conclusions, and those who disclose positions and let everyone challenge conclusions.

Why is that discipline vital? Consider an extreme failure. In the Soviet Union of the 1930s and 1940s, the agronomist Trofim Lysenko promoted a theory rejecting Mendelian genetics and claiming that environment could directly remake heredity. It rose not by surviving tests but by fitting the officially favored worldview. Dissenters lost posts or were exiled. Crop breeding took decades of wrong turns and agriculture paid heavily. Values here did not merely select questions; they rewrote what counted as evidence. Weber's discipline had been dismantled. The example shows that public evidence, permitted rebuttal, and conclusions separated from wishes are not scholarly fussiness. They are safeguards against burning the world with wishes.

We can now offer a decent answer to the title. Facts, explanations, and values \textbf{cannot be detached}, and we should not pretend otherwise. Values always help choose what we study and care about. But they \textbf{must be distinguished}, each thread identified and submitted to its appropriate tests. Values may determine questions, not falsify evidence. Explanations may be sympathetic, not absolving. Facts may be unpleasant, but owe nobody an apology.

\section{Three Baskets in the Trending-News Age}
The skill philosophers and historians have refined for millennia is now yours to use on a phone, because the internet makes a battle among the baskets an everyday sight.

First, reversal stories. A viral video shows an elder fallen by the road and a young person walking past. Comments instantly judge: people are coldhearted now, explanation; society is finished, value. Two days later, full footage reveals the young person had gone to get help. The initial comments failed not by caring about justice but by pouring a few seconds of footage from basket one straight into a final judgment on society in basket three, skipping the middle question of what happened and why. Stories reverse repeatedly not because facts are fickle, but because people rush to verdict before facts arrive.

Next, an upgraded version of position-first thinking. Online, people dismiss one another with ``Where you sit decides what you think,'' meaning a view reflects only one's interests. Is there something to it? Yes. Values affect what we notice and care about, and claims of positionlessness deserve suspicion. But as a universal weapon, it crosses a line: \textbf{every argument becomes merely a position in makeup}. Evidence and methods need no inspection; identify the speaker's camp and close the case. This is an everyday version of Lysenkoism, replacing whether a claim is true with who makes it. The healthy response is the opposite: because a person has a position, examine their evidence and reasoning \textbf{more carefully}, neither exempting nor dismissing them.

You may know the Chinese online label lizhongke, mocking a supposedly rational, impartial observer who really fudges the issue. It targets a real problem. Some people pretend equal blame is objectivity, posing neutrally between perpetrators and victims, itself a deeply nonneutral act. But the label can hit the wrong target. True objectivity is not positionlessness but declared positions and openness to testing. Someone can stand firmly with the vulnerable while rigorously checking every fact. These protect one another. A cause of justice is often damaged most not by enemies, but by an unverified piece of evidence offered by its own side.

Even conversations with artificial intelligence cannot escape the chapter. ``Give me a completely objective, neutral answer'' is itself a trap. Machine responses come from human texts containing countless choices about what to write, omit, or foreground. A machine's lack of a position does not mean its output lacks positions. It blends thousands of human positions and serves them wearing an opinionless mask. The tool is especially useful here. Separate factual claims, which can be checked; explanations, which compete with alternatives; and value judgments, asking whose side they take. Do this with machines, people, and yourself.

Finally, a school incident. You see your desk mate cheating. Should you tell the teacher? Sort the thoughts passing through your mind: they copied, fact, but did you see clearly? Family pressure caused it, explanation, perhaps true or not. Reporting a friend is wrong, value. Unpunished cheating is unfair to rule-followers, also value. Sorting does not make the decision easy or decide for you, but it reveals the level at which you hesitate. Some people argue all their lives without seeing that.

\section{Over to You: Dig or Do Not Dig?}
Return to the opening cave.

Suppose archaeologists find clues at a site suggesting that further excavation could disprove a local community's inherited origin story. The story tells children where they came from and why they live on this land; they sing it at festivals. Hearing the news, its keepers ask the team to stop. The facts you seek are part of our identity, they say. Unearth them and our song becomes impossible to sing.

The team divides. Some insist facts are facts, truth owes no apology: dig. Others argue knowledge is not free and its harms belong in the costs; at least negotiate a mutually acceptable plan with local people first. Others would let the story's keepers decide, provoking an immediate reply: what if people elsewhere use the same reason to ban research into slavery or war crimes?

Now the pen is yours. Dig or do not dig?

Give an answer and reasons, but first check your tools.

You know a river separates pollen from funeral, and conclusions cannot jump it for evidence. Values may choose research questions but should not alter what excavation finds. Objectivity means exposing methods, evidence, and reasoning to everyone, not having no position. You also know that sympathetically understanding why a community needs its song and writing that evidence points elsewhere can coexist in one heart.

Must truth be gentle? May feelings veto evidence? What price may a society pay for truth, and how much may it hide to protect people's feelings?

There is no standard answer. But whatever you write, you have used three kinds of sentence at once: stating observed facts, explaining others' thinking, declaring what should happen. The fine line between them is invisible and intangible, yet drawing it is the first step in learning honesty.
